\documentclass[11pt]{letter}
\usepackage[margin=1in]{geometry}
\usepackage{times}
\usepackage[hidelinks]{hyperref}

\signature{[Author Name]}
\address{Institution \\ Address}

\begin{document}

\begin{letter}{%
  Program Chairs \\
  CVPR 2026 EarthVision Workshop \\
}

\opening{Dear Program Chairs,}

This file is retained as a workshop-format cover-letter template for Paper 12. The active journal submission package is maintained under \texttt{submission/paper12\_isprs\_jprs\_20260606}; the wording below is synchronized with the current bounded manuscript rather than the earlier public-benchmark-only draft.

We are pleased to submit our manuscript, ``Architecture-Aware Diagnosis of Parameter-Efficient Adaptation in Prithvi-100M: Fused-QKV Failure, Capacity Boundaries, and Cross-Domain Remote-Sensing Validation,'' for consideration at the CVPR 2026 EarthVision Workshop.

\textbf{Summary.}
The paper studies how parameter-efficient fine-tuning (PEFT) methods behave on Prithvi-100M when sensor inputs, channel bridges, deployment labels, and transfer domains differ from pre-training. We evaluate seven adaptation strategies across EuroSAT, BigEarthNet-S2, LandCover.ai, a production-style Linhe County LULC workflow supervised by Esri-derived labels, and LoveDA urban/rural cross-domain transfer. The current manuscript also includes a completed EuroSAT capacity sweep, a channel-bridge ablation, and a Linhe synthetic weak-label control that is not an independent manual validation set. Three bounded findings emerge:

\begin{enumerate}
    \item Standard module-wise LoRA is vulnerable to a fused-QKV insertion trap on Prithvi-100M, while corrected split-QKV LoRA still remains below Houlsby adapters in the tested Prithvi-100M setting.
    \item Houlsby adapters provide the most reliable PEFT baseline in our experiments, but the result is scoped to one backbone and should not be read as a universal GeoFM ranking.
    \item Domain shift, label source, channel bridging, and adapter capacity materially affect the operating point, with LoveDA supporting an adapter-capacity hypothesis rather than a general rule for all geospatial foundation models.
\end{enumerate}

\textbf{Relevance to EarthVision.}
The paper addresses a practical gap in geospatial AI: practitioners often transfer PEFT intuitions from NLP and conventional vision without checking whether remote-sensing backbones expose the same trainable tensors or tolerate the same adaptation modules. Our benchmark and audit files make that inspection reproducible for Prithvi-100M, including the fused-attention insertion check, channel-bridge sensitivity, trainable-parameter accounting, and per-class cross-domain diagnostics.

\textbf{Novelty and contribution.}
The contribution is diagnostic rather than universal. We separate implementation failure from post-fix low-rank underperformance, show that input-stage adaptation is regime-dependent, and demonstrate that deployment-style labels and geographic scene splits can change the apparent PEFT delta. The result is a stress-test protocol for Prithvi-like GeoFM adaptation rather than a claim that any method replaces a full GIS production stack.

\textbf{Confirmation.}
This manuscript has not been published or submitted elsewhere in the form described here. All authors have approved the submission.

\closing{Sincerely,}

\end{letter}
\end{document}